% ==============================================================================
% CHAPTER 1: INTRODUCTION
% Chair of Wireless Systems (Fachgebiet Drahtlose Systeme)
% ==============================================================================

\chapter{Introduction}
\label{ch:introduction}

\begin{takeawaybox}[Writing Guide: Chapter 1 --- Introduction]
\textbf{Goal}: Hook the reader, establish the real-world and scientific context, define the precise research gap, formalize your research questions, and outline your core contributions.\\
\textbf{Typical Length}: 5--8 pages (approx. 8--10\% of the thesis).\\
\textbf{Key Components}:
\begin{enumerate}[leftmargin=*]
    \item \textbf{Motivation \& Context}: Why is this problem important now? (e.g. edge computing, autonomous wireless networks, LLM reliability).
    \item \textbf{Problem Statement \& Research Gap}: What is the specific limitation in existing literature or engineering practice?
    \item \textbf{Research Questions (RQs) \& Hypotheses}: 2 to 4 crisp, answerable research questions.
    \item \textbf{Proposed Approach}: A high-level overview of your solution before diving into details.
    \item \textbf{Key Contributions}: Bulleted list of tangible scientific and engineering deliverables.
    \item \textbf{Scope and Delimitations}: What is deliberately inside vs. outside the scope of this thesis.
    \item \textbf{Thesis Outline}: Roadmap guide to the subsequent chapters.
\end{enumerate}
\end{takeawaybox}

\section{Motivation and Context}
\label{sec:intro_motivation}

The rapid convergence of \acp{WSN}, edge computing architectures, and modern \ac{AI} has unlocked unprecedented capabilities in autonomous environmental monitoring, cyber-physical systems, and industrial IoT~\cite{buttyan2010application,piotrowski2006public}. Autonomous decision-making agents deployed at the network edge increasingly utilize foundation models, such as \acp{LLM}, to interpret noisy telemetry, orchestrate subtasks, and dynamically adapt to volatile wireless channel conditions~\cite{pechlivanidou2022survey,shi2016edge,mao2017survey}. As established by \textcite{meissner2026human}, value creation in modern intelligent systems is undergoing a fundamental paradigm shift from deterministic algorithmic execution (DevOps) toward probabilistic socio-technical orchestration (AgentOps) and joint cognitive systems.

However, operating machine learning pipelines in real-time, distributed edge environments introduces severe dependability challenges. Unlike centralized cloud clusters, edge nodes operate under stringent energy, memory, and communication bandwidth constraints~\cite{piotrowski2006public}. When an edge model encounters unexpected sensory data, ambiguous inputs, or out-of-distribution events, it must decide whether to process the query locally, escalate to a larger cloud model, or request assistance from a human operator~\cite{kendall2017uncertainties}.

\section{Problem Statement and Research Gap}
\label{sec:intro_problem}

Current orchestration dispatchers predominantly treat predictive uncertainty as a monolithic scalar quantity (e.g., softmax entropy or raw confidence scores). This simplistic treatment creates a severe systemic failure mode:
\begin{itemize}
    \item \textbf{Conflation of Uncertainty Types}: Standard dispatchers cannot distinguish between \emph{epistemic uncertainty} (which stems from a model's lack of knowledge and can be resolved by retrieving additional context or consulting a larger model) and \emph{aleatoric uncertainty} (which stems from inherent ambiguity, sensor noise, or irremediable stochasticity in the environment)~\cite{kuhn2023semantic,kendall2017uncertainties}.
    \item \textbf{Suboptimal Routing \& Resource Waste}: Treating all high-uncertainty events identically causes pipelines to waste computational energy and network bandwidth escalating aleatoric queries to high-tier cloud models, which fail just as reliably as small models because the ground truth itself is fundamentally ambiguous.
    \item \textbf{Excessive Human Oversight Costs}: Conversely, systems defer benign, solvable epistemic queries to costly human operators, overwhelming human oversight in large-scale sensor deployments.
\end{itemize}

\section{Research Questions and Hypotheses}
\label{sec:intro_rqs}

To address this gap, this thesis investigates the design, implementation, and empirical validation of an uncertainty-type-driven orchestration architecture for distributed wireless systems. We formulate the following core research questions:

\begin{researchquestion}[Research Question 1 (RQ1)]
\textbf{RQ1: Orthogonal Uncertainty Decomposition}\\
How can epistemic and aleatoric uncertainty signals be reliably and efficiently separated in sequential decision pipelines using lightweight semantic clustering at the wireless edge?
\end{researchquestion}

\begin{researchquestion}[Research Question 2 (RQ2)]
\textbf{RQ2: Routing Efficiency and Resource Optimization}\\
To what extent does routing subtasks conditioned on uncertainty type reduce unnecessary cloud escalations and inference latency compared to scalar confidence baselines?
\end{researchquestion}

\begin{researchquestion}[Research Question 3 (RQ3)]
\textbf{RQ3: Human-in-the-Loop Deferral Reliability}\\
How does type-driven selective prediction affect the precision of human-in-the-loop escalations in mission-critical sensor monitoring tasks?
\end{researchquestion}

\section{Proposed Approach: UTSR in Brief}
\label{sec:intro_approach}

This work introduces \ac{UTSR}, a modular routing framework designed for resilient edge intelligence. \ac{UTSR} decomposes prediction variance into orthogonal epistemic and aleatoric metrics by coupling semantic entropy estimation with bidirectional Natural Language Inference (\ac{NLI}). When epistemic uncertainty dominates, the subtask is routed to local vector memory retrieval (\ac{RAG}) or an upstream model. When aleatoric uncertainty dominates, the subtask triggers multi-hypothesis disambiguation or selective deferral to human supervisors, eliminating wasteful computation.

\section{Key Contributions}
\label{sec:intro_contributions}

The primary contributions of this Master's thesis are summarized as follows:
\begin{enumerate}
    \item \textbf{Formal Framework for Uncertainty Routing}: We formalize the distinction between epistemic and aleatoric dispatching in distributed edge pipelines.
    \item \textbf{Lightweight Edge-Compatible Implementation}: We develop a concrete implementation of \ac{UTSR} optimized for resource-constrained edge hardware.
    \item \textbf{Empirical Benchmark Evaluation}: We evaluate the approach across established multi-hop reasoning datasets (HotpotQA~\cite{yang2018hotpotqa}) and ambiguous QA datasets (AmbigQA~\cite{min2020ambigqa}).
    \item \textbf{Physical Sensor Network Validation}: We deploy \ac{UTSR} on an ultrasonic wireless sensor testbed, measuring real-world latency, energy consumption, and communication overhead.
\end{enumerate}

\section{Scope and Delimitations}
\label{sec:intro_scope}

This thesis focuses specifically on inference-time routing and decision orchestration for natural language and sensor-derived telemetry. Full model parameter fine-tuning and hardware ASIC acceleration are considered outside the direct scope of this study.

\section{Thesis Outline}
\label{sec:intro_outline}

The remainder of this thesis is structured as follows:
\begin{itemize}
    \item \Cref{ch:background} provides the theoretical foundations of wireless edge architectures, uncertainty quantification, and deep learning pipelines.
    \item \Cref{ch:related_work} surveys the state of the art in edge orchestration and uncertainty estimation, highlighting the research gap.
    \item \Cref{ch:methodology} formalizes the \ac{UTSR} architecture and mathematical formulation.
    \item \Cref{ch:implementation} details the software stack, edge testbed configuration, and dataset preprocessing.
    \item \Cref{ch:evaluation} presents the empirical results, benchmark performance, and ablation studies.
    \item \Cref{ch:discussion} discusses experimental findings, architectural trade-offs, and threats to validity.
    \item \Cref{ch:conclusion} concludes the thesis, answers the research questions, and outlines avenues for future research.
\end{itemize}
